\chapter{Files d'attente, tampons circulaires et mémoire partagée}\label{ch:ll:queues-ring-buffers-and-shared-memory}

En 2011, les ingénieurs d'une bourse de détail publièrent les chiffres qui les avaient conduits à repenser leur système. Faire transiter un
événement par un pipeline à trois étapes, construit sur la file bloquante standard de Java, prenait 32\,757 nanosecondes par saut
en moyenne, et son 99.99\up{e}~centile dépassait quatre millisecondes~; avec le \omterm{def:ll:queues-ring-buffers-and-shared-memory:ring}{tampon circulaire} qu'ils construisirent à la place, 52
nanosecondes en moyenne et sous 8\,192 au 99.99\up{e}~centile (Thompson et al., 2011). Le travail accompli dans chaque étape était le
même~; ce sont les files qui faisaient la latence. Un processus de trading est un pipeline de threads (feed handler, carnet, stratégie, passerelle,
journalisation), et la manière dont les données circulent entre eux détermine en grande partie sa queue de distribution. Ce chapitre construit les structures qui assurent ce transfert~:
l'anneau à un producteur et un consommateur, l'anneau de diffusion du motif «~disruptor~», le \omterm{def:ll:queues-ring-buffers-and-shared-memory:seqlock}{verrou à séquence} pour les données où seule
la dernière valeur compte, et ces mêmes structures entre processus en mémoire partagée, avec la conduite à tenir quand un
consommateur ne peut pas suivre.

\section{L'anneau à un producteur et un consommateur}

\begin{definition}[Tampon circulaire, file à un producteur et un consommateur]\label{def:ll:queues-ring-buffers-and-shared-memory:ring}
Un \emph{tampon circulaire}\index{tampon circulaire} est un tableau de cases de taille fixe, utilisé de manière circulaire~: la position d'écriture et la position de lecture
croissent indéfiniment et indexent le tableau modulo sa taille, qui est une puissance de deux afin que le modulo soit un masque. Une \emph{file à un producteur et un
consommateur}\index{file à un producteur et un consommateur} (SPSC) est un anneau ayant exactement un thread d'écriture et un
thread de lecture~; chaque position n'est écrite que par un seul thread, de sorte qu'aucun \omterm{def:ll:lock-free-programming:cas}{compare-and-swap} n'est nécessaire.
\end{definition}

Le producteur écrit la case située à sa position, puis publie la nouvelle position par une écriture à sémantique release~; le consommateur charge
la position du producteur avec une lecture acquire, lit toutes les cases qui la précèdent, et publie sa propre position de la même façon, en rendant les
cases (chapitre~11). Les deux opérations sont \omterm{def:ll:lock-free-programming:progress}{sans attente} (wait-free). Lamport démontra la correction d'une telle file \omterm{def:ll:lock-free-programming:progress}{sans verrou} en 1983~; ce que
des décennies de pratique y ont ajouté, ce sont les détails qui la rendent rapide. Les deux positions occupent des \omterm{def:ll:memory-hierarchy-and-caches:line}{lignes de cache} distinctes, afin que les
écritures du producteur n'invalident pas la position du consommateur (chapitre~3). Et chaque côté conserve une copie privée de la position de l'autre,
qu'il ne rafraîchit que lorsque l'anneau paraît plein ou vide~: tant qu'il n'est ni l'un ni l'autre, un message ne coûte qu'un transfert de \omterm{def:ll:memory-hierarchy-and-caches:line}{ligne de cache}
(la case) au lieu de trois.

\begin{omfigure}
\begin{tikzpicture}[font=\footnotesize]
\foreach \i/\c in {0/5,1/5,2/30,3/30,4/30,5/5,6/5,7/5} {
  \pgfmathsetmacro{\a}{90 - 45*\i}
  \draw[fill=omDef!\c,draw=omDef] (\a+22.5:1.3) arc (\a+22.5:\a-22.5:1.3) -- (\a-22.5:2.3) arc (\a-22.5:\a+22.5:2.3) -- cycle;
  \node at (\a:1.8) {\i};
}
\node[align=center] at (0,0) {8 cases\\ (masque 7)};
\draw[-Stealth,very thick,omThm] (3.9,0.9) node[right,align=left]{lecture = 2\\ (ligne du consommateur)} -- (0:2.4);
\draw[-Stealth,very thick,omMeth] (-3.9,-2.3) node[left,align=right]{écriture = 5\\ (ligne du producteur)} -- (-135:2.4);
\node[anchor=west,align=left] at (3.9,-1.0) {grisées~: cases 2, 3, 4,\\ écrites, pas encore lues};
\node[anchor=west,align=left] at (3.9,-1.9) {plein lorsque écriture $-$ lecture $=$ 8};
\end{tikzpicture}

\omcaption{Un anneau de huit cases après que cinq messages ont été écrits et deux lus (l'état de l'implantation partagée de la construction).
Les positions du producteur et du consommateur occupent chacune une \omterm{def:ll:memory-hierarchy-and-caches:line}{ligne de cache} qui lui est propre et ne font que croître~; la case de la position $p$ est
$p \bmod 8$.}\label{fig:ll:queues-ring-buffers-and-shared-memory:ring}
\end{omfigure}

La \Cref{fig:ll:queues-ring-buffers-and-shared-memory:rtt} mesure un aller-retour entre deux threads épinglés~: l'un envoie un
message de huit octets dans un anneau, l'autre le renvoie par un second anneau. Environ \qty{200}{\nano\second} à la
médiane~: quatre transferts de \omterm{def:ll:memory-hierarchy-and-caches:line}{ligne de cache} entre cœurs. Le même aller-retour par une file protégée par un mutex et une variable de condition
coûte environ \qty{30}{\micro\second}, soit plus de cent fois davantage, car le thread en attente s'endort dans
le noyau et doit être réveillé. Dans un seul sens, l'anneau débite environ 250 millions de messages de huit octets par seconde sur cet ordinateur portable.

\begin{omfigure}
\begin{tikzpicture}
\begin{axis}[width=12cm,height=5.8cm,ybar,bar width=7pt,ymode=log,log origin y=infty,ymin=100,ymax=5000000,
  symbolic x coords={threads,other-pair,processes,locked},
  xtick=data,xticklabels={{anneau SPSC,\\ threads},{anneau SPSC,\\ autre paire de CPU},{anneau SPSC,\\ deux processus},{mutex et\\ variable de condition}},
  xticklabel style={align=center,font=\footnotesize},enlarge x limits=0.15,
  ylabel={ns par aller-retour (échelle log)},tick label style={font=\small},label style={font=\small},grid=major,
  grid style={gray!25},legend columns=4,legend style={font=\footnotesize,at={(0.5,-0.33)},anchor=north,draw=none,fill=none,
  /tikz/every even column/.append style={column sep=6pt}},table/col sep=comma]
\addplot[fill=omDef!70,draw=omDef] table[x=key,y=p50]{figdata/low-latency/12-queues-ring-buffers-and-shared-memory/measured_rtt.csv};
\addplot[fill=omDef!35,draw=omDef] table[x=key,y=p99]{figdata/low-latency/12-queues-ring-buffers-and-shared-memory/measured_rtt.csv};
\addplot[fill=omThm!45,draw=omThm] table[x=key,y=p999]{figdata/low-latency/12-queues-ring-buffers-and-shared-memory/measured_rtt.csv};
\addplot[fill=gray!40,draw=gray] table[x=key,y=max]{figdata/low-latency/12-queues-ring-buffers-and-shared-memory/measured_rtt.csv};
\legend{p50,p99,p99.9,maximum}
\end{axis}
\end{tikzpicture}

\omcaption{Allers-retours d'un message de huit octets entre deux threads (ou processus) épinglés, par deux anneaux SPSC, et par
deux files protégées par un mutex et une variable de condition~; 100\,000 allers-retours chacun. Mesuré sur un ordinateur portable (Intel Core Ultra 7
155H) sous WSL2, sans cœurs isolés. Données~: \texttt{bench\_rings.py}.}\label{fig:ll:queues-ring-buffers-and-shared-memory:rtt}
\end{omfigure}

\section{Plusieurs lecteurs~: le motif disruptor}

\begin{definition}[Motif disruptor]\label{def:ll:queues-ring-buffers-and-shared-memory:disruptor}
Dans le \emph{motif disruptor}\index{motif disruptor}, un seul anneau transporte chaque événement vers plusieurs consommateurs, chacun lisant à sa
propre position~; les consommateurs qui dépendent les uns des autres attendent les positions des autres plutôt que des files placées entre eux, si bien que
un pipeline d'étapes devient un seul anneau à plusieurs curseurs, sans copies.
\end{definition}

Les données de marché en sont le cas naturel~: un feed handler, plusieurs consommateurs (le constructeur de carnet, une seconde stratégie, l'enregistreur,
un moniteur de risque), dont chacun a besoin de chaque message. Deux politiques diffèrent quant à celui qui attend. Avec la \omterm{def:ll:queues-ring-buffers-and-shared-memory:bp}{contre-pression} (back-pressure, section suivante), le
producteur attend le consommateur le plus lent, ce qui convient à une voie d'ordres, où rien ne doit être perdu. Pour les données de marché, le
producteur ne doit pas attendre~: la bourse ne ralentira pas. L'anneau de diffusion de la construction ne bloque donc jamais son producteur~; chaque
case porte le numéro de séquence du message qu'elle contient, de sorte qu'un lecteur qui a été dépassé voit une séquence plus récente que celle
qu'il attendait et sait qu'il a perdu des messages, au lieu de lire un mélange.

\begin{omfigure}
\begin{tikzpicture}[font=\footnotesize]
\foreach \i/\c in {0/8,1/9,2/10,3/11,4/12,5/13,6/14,7/15} {
  \draw[fill=omDef!10,draw=omDef] (1.3*\i,0) rectangle ++(1.2,0.8);
  \node at (1.3*\i+0.6,0.52) {case \i};
  \node[font=\scriptsize] at (1.3*\i+0.6,0.2) {msg \c};
}
\draw[-Stealth,very thick,omMeth] (0.6,1.6) node[above]{producteur~: suivant = 16, dans la case 0} -- (0.6,0.85);
\draw[-Stealth,thick,omDef] (7.1,-0.8) -- (7.1,-0.05);
\node[anchor=north east,align=right,omDef] at (7.0,-0.75) {enregistreur\\ à 13};
\draw[-Stealth,thick,omDef] (9.7,-0.8) -- (9.7,-0.05);
\node[anchor=north west,align=left,omDef] at (9.8,-0.75) {constructeur de carnet\\ à 15};
\draw[-Stealth,thick,omThm] (8.4,-1.95) -- (8.4,-0.05);
\node[anchor=north,align=center,omThm] at (8.4,-1.95) {moniteur lent à 6~: case 6\\ contient msg 14, dépassé};
\end{tikzpicture}

\omcaption{Un anneau de diffusion de huit cases après 16 messages. Chaque lecteur conserve sa propre position~; le producteur n'attend jamais. Un
lecteur à la position 6 trouve dans la case 6 le message 14, dont la séquence est plus récente que celle attendue~: il a été dépassé (\omterm{def:ll:queues-ring-buffers-and-shared-memory:bp}{consommateur lent})
et doit se resynchroniser.}\label{fig:ll:queues-ring-buffers-and-shared-memory:broadcast}
\end{omfigure}

\section{Verrous à séquence pour les données de dernière valeur}

\begin{definition}[Verrou à séquence]\label{def:ll:queues-ring-buffers-and-shared-memory:seqlock}
Un \emph{verrou à séquence}\index{verrou à séquence} (sequence lock) publie une petite valeur qu'un seul écrivain met à jour et que de nombreux lecteurs consultent~: l'écrivain
incrémente un compteur jusqu'à une valeur impaire, écrit la valeur, puis l'incrémente à nouveau jusqu'à une valeur paire~; un lecteur lit le compteur, copie la valeur,
relit le compteur, et recommence si les deux lectures diffèrent ou sont impaires.
\end{definition}

La documentation du noyau Linux décrit ce mécanisme comme reposant sur des lecteurs \omterm{def:ll:lock-free-programming:progress}{sans verrou}, avec des boucles de nouvelle tentative en lecture seule
et sans famine de l'écrivain, pour des données écrites rarement par rapport à leur fréquence de lecture. En trading, ce motif convient au haut du carnet
ou à un prix de référence, où les lecteurs veulent la dernière valeur cohérente, et non chacune des valeurs intermédiaires. Deux détails le rendent correct
en C++ comme sur x86~: la valeur doit être copiée au moyen d'\omterm{def:ll:lock-free-programming:atomic}{opérations atomiques} (les opérations relâchées suffisent), car une copie en concurrence avec
les écritures non atomiques de l'écrivain est une \omterm{def:ll:rust-for-low-latency:race}{course aux données} et un \omterm{def:ll:cpp-for-latency-iii-the-compiler:ub}{comportement indéfini} (Boehm, 2012)~; et la seconde lecture du compteur
du lecteur doit être ordonnée après la copie par une barrière d'acquisition. Le test de la construction fait tourner un écrivain dans une boucle serrée
face à un lecteur qui vérifie deux cent mille paires pour détecter les déchirures.

\section{Mémoire partagée entre processus}

\begin{definition}[Transport par mémoire partagée]\label{def:ll:queues-ring-buffers-and-shared-memory:shm}
Un \emph{transport par mémoire partagée}\index{transport par mémoire partagée} (shared-memory transport) achemine des messages entre processus par l'intermédiaire d'une région de mémoire
que chacun projette dans son espace d'adressage (sous Linux, un objet de mémoire partagée POSIX issu de \texttt{shm\_open}, ou un fichier sur un
système de fichiers en mémoire), organisée en anneaux dont les positions sont des atomiques~; aucun appel système n'est effectué par message.
\end{definition}

Les processus isolent les défaillances~: une stratégie qui plante n'entraîne pas nécessairement le feed handler avec elle, et des composants écrits dans
des langages différents peuvent partager un même flux. L'anneau s'en moque~: la \cref{fig:ll:queues-ring-buffers-and-shared-memory:rtt} montre
le même aller-retour entre deux processus que entre deux threads. Ce qui change, c'est le contrat entre les processus~: la disposition doit être spécifiée
octet par octet et versionnée, puisque deux binaires compilés à des dates différentes la lisent~; chaque champ doit avoir une taille et un alignement fixes~;
et rien dans la région ne peut être un pointeur, ce qui signifie que rien n'y est valable dans l'autre processus.
La disposition retenue par la construction est un en-tête de 192~octets (nombre magique, version, taille des cases, capacité, puis les deux positions,
chacune sur sa propre ligne), suivi des cases~; ses implémentations en C++ et en Rust reproduisent la même image binaire d'un anneau, écrites
indépendamment à partir de la spécification par un script.

\section{Contre-pression, consommateurs lents et conflation}

\begin{definition}[Contre-pression, consommateur lent, conflation]\label{def:ll:queues-ring-buffers-and-shared-memory:bp}
La \emph{contre-pression}\index{contre-pression} (back-pressure) est la propagation à son producteur de la lenteur d'un consommateur, lequel doit attendre ou
refuser de nouveaux travaux lorsque la file qui les sépare est pleine. Un \emph{consommateur lent}\index{consommateur lent} (slow consumer) est un consommateur dont le débit
est inférieur à celui de son entrée, de sorte que sa file s'allonge ou, sur un anneau qui n'attend pas, qu'il est dépassé. La \emph{conflation}\index{conflation}
est la livraison délibérée de la seule dernière valeur de chaque élément (un niveau de prix, une cotation) à un consommateur qui ne peut pas absorber
chaque mise à jour.
\end{definition}

La taille d'une file est un énoncé sur les rafales. Le livre~4, chapitre~8, donne la réponse de la file M/M/1 pour un étage dont les arrivées suivent un processus
de Poisson de taux $\lambda$ et dont le service est exponentiel de taux $\mu$~: avec un taux d'utilisation $\rho = \lambda/\mu$, le nombre stationnaire
dans le système dépasse $n$ avec une probabilité $\rho^n$. Un anneau qui doit déborder moins d'une fois sur un milliard de messages à $\rho = 0.9$
a donc besoin de $n$ tel que $0.9^n < 10^{-9}$, soit environ 197 cases, 256 en puissance de deux. Les données de marché sont plus rafaleuses qu'un processus
de Poisson~: le dimensionnement part de la valeur de la file M/M/1 et se vérifie sur des rafales enregistrées (chapitre~18).

\begin{proposition}[Temps avant qu'un consommateur lent soit dépassé]\label{prop:ll:queues-ring-buffers-and-shared-memory:lap}
Un anneau de diffusion de $N$ cases dont le producteur écrit au taux $\lambda_b$ pendant une rafale, lu par un consommateur au taux $\mu < \lambda_b$
qui démarre la rafale avec un retard nul, dépasse le consommateur au bout de $N/(\lambda_b - \mu)$ secondes.
\end{proposition}

\begin{proof}
Le retard du consommateur croît de $\lambda_b - \mu$ messages par seconde~; le producteur écrase la prochaine case non lue du consommateur lorsque
le retard atteint $N$.
\end{proof}

Un anneau de 65\,536 cases, une rafale à 2 millions de messages par seconde et un moniteur qui lit 1.5 million donnent au moniteur 131
millisecondes~; une rafale d'ouverture plus longue le dépasse. Les réponses de conception consistent à offrir à de tels consommateurs des vues
conflatées (un \omterm{def:ll:queues-ring-buffers-and-shared-memory:seqlock}{verrou à séquence} par instrument au lieu de chaque message), à dimensionner leurs anneaux sur des rafales mesurées, ou à les laisser
prendre du retard par choix et à se resynchroniser depuis un instantané, comme le fait le feed handler du chapitre~18 auprès de la bourse.

\section{Tutoriel~: anneaux dans les threads, les processus et deux langages}\label{tut:ll:queues-ring-buffers-and-shared-memory:tut}

\begin{tutorial}
\textbf{Objectif.} Construire les anneaux, vérifier la disposition partagée en C++ et en Rust sur une même référence, et mesurer les aller-retours
et le débit. \textbf{État final~:} \cref{fig:ll:queues-ring-buffers-and-shared-memory:rtt}, une figure de débit, et des tests au vert
dans les deux langages.
\begin{tutsteps}
\item \textbf{Écrire et publier.} Le producteur copie le message dans sa case et publie sa position par une écriture release~;
il ne lit la position du consommateur que lorsque l'anneau semble plein.
\omcode{code/firm/ring/cpp/firm_ring.hpp}{68}{95}{Le producteur et le consommateur de l'anneau SPSC.}\label{lst:ll:queues-ring-buffers-and-shared-memory:spsc}
\item \textbf{La même chose en Rust}, sur les mêmes octets~: une région de mots atomiques, écrite par l'intermédiaire de leur mutabilité intérieure.
\omcode{code/firm/ring/rust/src/lib.rs}{80}{101}{Le producteur Rust, qui écrit la disposition attendue par le lecteur C++.}\label{lst:ll:queues-ring-buffers-and-shared-memory:rust}
\item \textbf{Un \omterm{def:ll:queues-ring-buffers-and-shared-memory:seqlock}{verrou à séquence}} dont la valeur est copiée comme des mots atomiques relâchés.
\omcode{code/firm/ring/cpp/firm_ring.hpp}{155}{179}{Écrivain et lecteur du verrou à séquence.}\label{lst:ll:queues-ring-buffers-and-shared-memory:seqlock}
\item \textbf{Vérifier la disposition.} \texttt{make\_ring\_fixture.py} écrit l'image binaire d'un anneau à partir de la spécification~; les deux
suites de tests construisent le même anneau et comparent, puis consomment la référence.
\item \textbf{Mesurer} avec \texttt{python bench\_rings.py}~: aller-retours entre deux threads sur deux paires de CPU, entre deux
processus, et par une file verrouillée, puis débit à sens unique.
\end{tutsteps}
\textbf{Pour aller plus loin.} Supprimer les copies en cache de la position de l'autre côté et mesurer à nouveau le débit~; placer les deux positions sur
une même \omterm{def:ll:memory-hierarchy-and-caches:line}{ligne de cache} et mesurer l'aller-retour.
\end{tutorial}

\section{Construction~: les anneaux de la firme}\label{bld:ll:queues-ring-buffers-and-shared-memory:ring}

\begin{build}
\bfield{Rôle} Le transport de chaque composant de la partie~IV~: le feed handler publie des événements normalisés sur un anneau de diffusion
(chapitre~18), le moteur de stratégie les lit et transmet les ordres à la passerelle par un anneau SPSC (chapitres~20--21), le journaliseur binaire
vide un anneau (chapitre~23), et le haut de chaque carnet est publié par un \omterm{def:ll:queues-ring-buffers-and-shared-memory:seqlock}{verrou à séquence}.
\bfield{Interface} C++20 \texttt{firm::ring}~: \texttt{format(base, capacity, slot\_size)}, \texttt{region\_size},
\texttt{Spsc(base)} avec \texttt{try\_write(msg, len)} et \texttt{try\_read(out, cap)}~; \texttt{Broadcast(base)} avec
\texttt{write}, \texttt{read(pos, out, len)} renvoyant \texttt{Ok}, \texttt{Empty} ou \texttt{Overrun}, et \texttt{head()}~;
\texttt{SeqLock<T>} avec \texttt{store} et \texttt{load}~; \texttt{Segment(name, bytes, create)}. Rust \texttt{firm\_ring}~:
\texttt{Region}, \texttt{format}, \texttt{Spsc::attach} avec \texttt{try\_write} et \texttt{try\_read}, \texttt{SeqLock2}.
\bfield{Règles} La disposition d'octets ci-dessus, en petit-boutiste, versionnée par un nombre magique et une version~; les capacités sont des puissances
de deux~; les positions sur des \omterm{def:ll:memory-hierarchy-and-caches:line}{lignes de cache} distinctes~; l'anneau SPSC signale le plein et le vide, l'anneau de diffusion n'attend jamais et signale
les dépassements~; aucune allocation après la construction.
\bfield{Tests d'acceptation} \texttt{code/firm/ring/}~: les octets de la référence reproduits et consommés en C++ et en Rust~; 200\,000
messages dans l'ordre entre deux threads dans les deux langages~; détection des dépassements de l'anneau de diffusion et resynchronisation~; aucune lecture déchirée sur 200\,000
lectures de \omterm{def:ll:queues-ring-buffers-and-shared-memory:seqlock}{verrou à séquence} sous un écrivain dans une boucle serrée~; 1\,000 messages venant d'un processus enfant par mémoire partagée.
\bfield{Pour aller plus loin} Un anneau de diffusion dont la politique déclarée pour le consommateur le plus lent est choisie (attendre, abandonner, conflater)~;
un anneau dans un fichier sur un système de fichiers à \omterm{def:ll:memory-hierarchy-and-caches:tlb}{grandes pages}.
\end{build}

\begin{omsources}
\item M.~Thompson, D.~Farley, M.~Barker, P.~Gee et A.~Stewart, \emph{Disruptor: high performance alternative to bounded queues
for exchanging data between concurrent threads}, 2011.
\item L.~Lamport, ``Specifying concurrent program modules'', \emph{ACM TOPLAS} 5(2), 1983.
\item H.-J.~Boehm, ``Can seqlocks get along with programming language memory models?'', MSPC 2012~; documentation du noyau Linux,
``Sequence counters and sequential locks''.
\item Page de manuel Linux \texttt{shm\_open(3)}.
\end{omsources}

\section{Exercices}

\begin{exercise}[$\star$]\label{exo:ll:queues-ring-buffers-and-shared-memory:1}
Un anneau compte 1\,024 cases~; le producteur est à la position 70\,001 et le consommateur à 69\,500. Combien de messages attendent, quelle case
le producteur écrira-t-il ensuite, et combien peut-il encore écrire avant que l'anneau soit plein~?
\end{exercise}

\begin{exercise}[$\star$]\label{exo:ll:queues-ring-buffers-and-shared-memory:2}
À 250 millions de messages par seconde, combien de temps une position sur 64~bits met-elle à reboucler~? Et une position sur 32~bits~?
\end{exercise}

\begin{exercise}[$\star$]\label{exo:ll:queues-ring-buffers-and-shared-memory:3}
Un étage fonctionne à $\rho = 0.8$. Quelle taille d'anneau rend le débordement d'une file M/M/1 plus rare qu'une fois sur un million~?
\end{exercise}

\begin{exercise}[$\star\star$]\label{exo:ll:queues-ring-buffers-and-shared-memory:4}
Pourquoi chaque côté de l'anneau SPSC garde-t-il une copie privée de la position de l'autre~? Que coûte par message la suppression de ces copies~?
\end{exercise}

\begin{exercise}[$\star\star$]\label{exo:ll:queues-ring-buffers-and-shared-memory:5}
Un moniteur de risque lit l'anneau de diffusion à 400\,000 messages par seconde. Une rafale de 3 millions de messages par seconde dure 80
millisecondes sur un anneau de 131\,072 cases. Le moniteur est-il rattrapé par le producteur, et si oui, quand~?
\end{exercise}

\begin{exercise}[$\star\star$]\label{exo:ll:queues-ring-buffers-and-shared-memory:6}
Pourquoi la valeur du \omterm{def:ll:queues-ring-buffers-and-shared-memory:seqlock}{verrou à séquence} doit-elle être copiée avec des \omterm{def:ll:lock-free-programming:atomic}{opérations atomiques} en C++, alors que le lecteur écarte toute copie prise pendant une
écriture~?
\end{exercise}

\begin{exercise}[$\star\star\star$]\label{exo:ll:queues-ring-buffers-and-shared-memory:7}
\emph{Programmation.} Placer les positions du producteur et du consommateur sur la même \omterm{def:ll:memory-hierarchy-and-caches:line}{ligne de cache} (modifier les décalages d'en-tête dans une copie de
\texttt{firm\_ring.hpp}) et relancer les bancs d'essai d'aller-retour et de débit. Qu'est-ce qui change, et pourquoi~?
\end{exercise}

\begin{exercise}[$\star\star\star$]\label{exo:ll:queues-ring-buffers-and-shared-memory:8}
\emph{Trouver la faille.} «~Notre anneau en mémoire partagée stocke, dans son en-tête, un pointeur vers le message de tête et une \texttt{std::string} pour le
nom de l'instrument.~»
\end{exercise}

\section{Problème~: Trois stratégies, un seul flux}

\begin{problem}[{Problème du week-end --- dimensionner un anneau de diffusion}]\label{pb:ll:queues-ring-buffers-and-shared-memory:1}
Un feed handler publie des événements normalisés sur un anneau de diffusion de 65\,536 cases de 64 octets. Trois consommateurs le lisent~: un constructeur de carnet
(book builder) à 4 millions d'événements par seconde au plus, une stratégie à 2.5 millions et un moniteur à 1.5 million. Un jour ordinaire, le flux
fait en moyenne 300\,000 événements par seconde~; à l'ouverture, les rafales atteignent 2 millions par seconde pendant une demi-seconde au plus.

\medskip\noindent\textbf{Partie I --- L'anneau.}
\begin{enumerate}
\item Quelle place occupe l'anneau en octets, et tient-il dans le cache de deuxième niveau du chapitre~3~?
\item Quel est le taux d'utilisation de chaque consommateur un jour ordinaire~?
\item Quels consommateurs peuvent accumuler du retard pendant une rafale~?
\item Pourquoi le producteur n'attend-il jamais après eux~?
\end{enumerate}

\medskip\noindent\textbf{Partie II --- La rafale.}
\begin{enumerate}[resume]
\item Au bout de combien de temps le moniteur est-il rattrapé (\cref{prop:ll:queues-ring-buffers-and-shared-memory:lap})~?
\item Est-il rattrapé pendant une rafale d'une demi-seconde~?
\item Quelle taille d'anneau permettrait au moniteur de traverser la rafale~?
\item Combien de mémoire cela coûterait-il~?
\end{enumerate}

\medskip\noindent\textbf{Partie III --- Les solutions de rechange.}
\begin{enumerate}[resume]
\item Le moniteur n'a besoin que du dernier prix de chacun des 5\,000 instruments. Que coûte en mémoire un \omterm{def:ll:queues-ring-buffers-and-shared-memory:seqlock}{verrou à séquence} par instrument,
et quelle part des mises à jour le moniteur saute-t-il pendant la rafale~?
\item Que fait le moniteur lorsqu'il détecte un dépassement sur l'anneau~?
\item Pourquoi la \omterm{def:ll:queues-ring-buffers-and-shared-memory:bp}{contre-pression} (back-pressure) est-elle fausse pour ce producteur et juste pour la voie des ordres~?
\item Que dit le modèle M/M/1 de l'anneau de la stratégie à 300\,000 événements par seconde, et pourquoi est-il optimiste~?
\end{enumerate}

\medskip\noindent\textbf{Partie IV --- Le verdict.}
\begin{enumerate}[resume]
\item Énoncer le \emph{résultat nommé}~: le temps avant que le moniteur soit rattrapé pendant une rafale, et la part des mises à jour qu'un moniteur à
\omterm{def:ll:queues-ring-buffers-and-shared-memory:bp}{conflation} saute.
\item Que dit l'aller-retour mesuré du coût d'ajout d'un consommateur~?
\item Pourquoi la disposition de l'anneau est-elle versionnée~?
\item Que prédirait ici, pour une file bloquante, la comparaison publiée entre une file bloquante et un anneau~?
\item Quelle règle de placement du chapitre~4 s'applique au feed handler et au constructeur de carnet~?
\item Comment enregistrer les tailles de rafale que l'anneau doit porter~?
\item Que doit faire la stratégie après un trou dans son entrée~?
\item En une phrase~: la taille d'un \omterm{def:ll:queues-ring-buffers-and-shared-memory:ring}{tampon circulaire} est une affirmation sur quoi~?
\end{enumerate}
\end{problem}

\section{Questions d'entretien}

\begin{interviewq}[$\star$ \iqroles{développeur}]\label{iq:ll:queues-ring-buffers-and-shared-memory:1}
Implémentez un \omterm{def:ll:queues-ring-buffers-and-shared-memory:ring}{tampon circulaire} à un producteur et un consommateur (SPSC). Quels \omterm{def:ll:lock-free-programming:ordering}{ordonnancements mémoire} lui faut-il~?
\end{interviewq}

\begin{interviewq}[$\star\star$ \iqroles{développeur}]\label{iq:ll:queues-ring-buffers-and-shared-memory:2}
Pourquoi un \omterm{def:ll:queues-ring-buffers-and-shared-memory:ring}{tampon circulaire} est-il plus rapide qu'une file protégée par un mutex et une variable de condition~?
\end{interviewq}

\begin{interviewq}[$\star\star$ \iqroles{développeur}]\label{iq:ll:queues-ring-buffers-and-shared-memory:3}
Qu'est-ce qu'un \omterm{def:ll:queues-ring-buffers-and-shared-memory:seqlock}{verrou à séquence}, quand l'utiliseriez-vous, et quels sont ses pièges~?
\end{interviewq}

\begin{interviewq}[$\star\star$ \iqroles{développeur}]\label{iq:ll:queues-ring-buffers-and-shared-memory:4}
Un flux de données de marché, trois consommateurs de vitesses différentes. Concevez la distribution.
\end{interviewq}

\begin{interviewq}[$\star\star$ \iqroles{développeur}]\label{iq:ll:queues-ring-buffers-and-shared-memory:5}
Que doit éviter une structure de données partagée en mémoire entre deux processus~?
\end{interviewq}

\begin{interviewq}[$\star\star\star$ \iqroles{développeur, chercheur}]\label{iq:ll:queues-ring-buffers-and-shared-memory:6}
Comment dimensionneriez-vous la file placée devant un étage dont l'entrée est en rafales~? Qu'apporte la théorie des files d'attente, et ce
qu'elle laisse échapper~?
\end{interviewq}
